\chapter{Manual de usuario}
\label{app:manual_usuario}

\ph{Este apéndice está dirigido al usuario final del sistema. A diferencia
del manual de instalación, no asume conocimientos técnicos y se centra en
\emph{cómo usar} el sistema una vez en marcha. Incluye capturas de pantalla
de las acciones principales cuando completes la memoria final.}

\section{Acceso al sistema}
\label{sec:acceso}

\ph{Indica la URL (o forma de acceso) del sistema y los pasos para
registrarse e iniciar sesión.}

\begin{itemize}
  \item \ph{URL de acceso: \url{https://example.org}}
  \item \ph{Registro de cuenta nueva.}
  \item \ph{Inicio de sesión.}
  \item \ph{Recuperación de contraseña (si aplica).}
\end{itemize}

\ph{(Espacio reservado para una captura de la pantalla de inicio.)}

\section{Funcionalidad principal}
\label{sec:funcionalidad_principal}

\ph{Describe paso a paso cómo realizar la tarea principal del sistema.
Conviene incluir capturas o anotaciones que destaquen los botones y campos
relevantes.}

\begin{enumerate}
  \item \ph{Paso 1: dónde hacer clic, qué introducir.}
  \item \ph{Paso 2: qué pantalla aparece y qué hacer en ella.}
  \item \ph{Paso 3: cómo confirmar la acción.}
  \item \ph{Paso 4: cómo verificar el resultado.}
\end{enumerate}

\section{Funcionalidades secundarias}
\label{sec:funcionalidades_secundarias}

\ph{Describe el resto de funcionalidades relevantes para el usuario, una
por subsección.}

\subsection{\ph{Gestión de \dots}}
\ph{Cómo crear, listar, editar y eliminar el recurso correspondiente.}

\subsection{\ph{Consulta de \dots}}
\ph{Cómo acceder a la información, qué filtros y ordenaciones están
disponibles.}

\subsection{\ph{Gestión del perfil de usuario}}
\ph{Cómo modificar datos personales, cambiar la contraseña y cerrar sesión.}

\section{Mensajes de error y resolución de incidencias}
\label{sec:errores}

\ph{Lista de errores frecuentes que puede encontrar el usuario, con su
significado y la acción recomendada para resolverlos.}

\begin{itemize}
  \item \textbf{\ph{Mensaje 1}}: \ph{causa habitual y solución.}
  \item \textbf{\ph{Mensaje 2}}: \ph{causa habitual y solución.}
  \item \textbf{\ph{Mensaje 3}}: \ph{causa habitual y solución.}
\end{itemize}

\section{Contacto y soporte}
\label{sec:contacto}

\ph{Vía de contacto para incidencias o soporte (correo, formulario,
repositorio de issues), si la hubiese.}
